% Part: computability
% Chapter: computability-theory
% Section: coding-computations

\documentclass[../../../include/open-logic-section]{subfiles}

\begin{document}

\olfileid{cmp}{thy}{cod}
\olsection{સંગણનોનું સંકેતીકરણ}

સંગણનાના દરેક નિદર્શમાં નીચેનું કરી શકાય છે:
\begin{enumerate}
\item સંગણનીય વિધેયોની \emph{વ્યાખ્યાઓ} વ્યવસ્થિત રીતે વર્ણવી શકાય.
  દાખલા તરીકે, ટ્યુરિંગ મશીનનાં વિનિર્દેશો, પુનરાવર્તી વ્યાખ્યાઓ
  અથવા કોઈ પ્રોગ્રામિંગ ભાષાના પ્રોગ્રામો આવી વ્યાખ્યાઓ આપે છે.
\item કોઈ આપેલી વ્યાખ્યા અને ઇનપુટ માટે વિધેયની સંગણનાનો પૂરો
  ઇતિહાસ વર્ણવી શકાય. દાખલા તરીકે, ટ્યુરિંગ મશીનની સંગણનાને
  તેના દરેક તબક્કાની ગોઠવણીના અનુક્રમથી વર્ણવી શકાય છે; તેમાં
  મશીનની સ્થિતિ અને ટેપ પરનું લખાણ આવે છે.
\item સંગણનાનો દેખીતી ઇતિહાસ ખરેખર આપેલી વ્યાખ્યાવાળું સંગણનીય
  વિધેય આપેલા ઇનપુટ પર કેવી રીતે ગણાય છે તેનો જ ઇતિહાસ છે કે
  નહીં તે તપાસી શકાય; અહીં એ ઇનપુટ પર વિધેય વ્યાખ્યાયિત છે,
  એટલે સંગણના અટકે છે.
\item સંગણનાના પૂરા ઇતિહાસના આવા વર્ણનમાંથી આપેલા ઇનપુટ માટે
  વિધેયનું મૂલ્ય કાઢી શકાય. દાખલા તરીકે, અટકતી ટ્યુરિંગ મશીનની
  સંગણનાના છેલ્લા તબક્કે ટેપ પરનું લખાણ એ મૂલ્ય છે.
\end{enumerate}

સંકેતીકરણથી સંગણનીય વિધેયના દરેક વર્ણનને એક સંખ્યાત્મક
\emph{સૂચક} આપી શકાય છે, એવી રીતે કે એ સૂચકમાંથી સૂચનાઓ
સંગણનીય રીતે ફરી મેળવી શકાય. એ જ રીતે સંગણનાનો પૂરો ઇતિહાસ
એક સંખ્યાથી સંકેતિત કરી શકાય. પરિણામે, ``$s$ એ ઇનપુટ~$x$
માટે સૂચક~$e$ ધરાવતા વિધેયની સંગણનાનો ઇતિહાસ સંકેતિત કરે છે''
એ અંકગણિતીય સંબંધ અને ``સંકેત~$s$ ધરાવતા સંગણના-અનુક્રમનું
પરિણામ'' એ વિધેય બંને સંગણનીય છે; ખરેખર, બંને આદિમ પુનરાવર્તી છે.

આ પાયાનો તથ્ય બહુ શક્તિશાળી છે. તે પસંદ કરેલા સંગણન-નિદર્શનથી
સ્વતંત્ર રીતે સંગણનીયતા વિશેના અનેક આશ્ચર્યકારક અને મહત્વના
પરિણામો સાબિત કરવા દે છે.

\end{document}
